\section{Pipeline Design}
% The template has no separate exploration section, so the data investigation
% (20% of the marks) is described here as the first stage of the pipeline.
% Sources: notebook cells 3-8 (exploration) and 13-14, 16, 20, 24 (model).

\begin{figure}[!htb]
\centering
\includegraphics[width=\textwidth]{pipeline.png}
\caption{Pipeline from raw data to the final evaluation. Every step under the bracket only sees training rows.}
\label{fig:pipeline}
\end{figure}

\textbf{Data investigation.}
BMI is missing for 201 patients (3.9\%), but 40 of them had a stroke, which is 16.06\% of all stroke cases, so dropping these rows would remove too much of the minority class.
The missing values were instead filled with the median BMI rather than the mean (28.10 and 28.89 on the full data; the model learns the median from the training rows only), because BMI is positively skewed (skew 1.06, with values up to 97.6; Figure~\ref{fig:bmi}).
The ID column was dropped, as was the single `Other' gender row, which had no stroke cases and so carries no information for stroke detection.
`Unknown' smoking status (1,544 patients, 30\%) was kept as its own category: it is not random, as these patients are much younger (mean age 30) and many are children, so filling it in would invent smoking histories for children.
Range checks found no impossible values: the binary columns contain only 0 and 1, and the 115 non-integer ages (at most 1.88) are infants recorded in months, so they are not errors.
Stroke rate rises sharply with age, from 0.1\% in the 18--33 band to 17.4\% in the 68--89 band (UKHSA age bands), and stroke patients also show a second peak at high glucose levels (Figure~\ref{fig:dist}).

\begin{figure}[!h]
\centering
\includegraphics[width=0.72\textwidth]{distributions.png}
\caption{Distributions of age, average glucose level and BMI for patients without (0) and with (1) a stroke.}
\label{fig:dist}
\end{figure}

Stroke rates are also much higher with hypertension (about 13\% vs 4\%) and heart disease (about 17\% vs 4\%), while ever-married patients (6.6\% vs 1.7\%) and children (almost 0\%) mostly reflect age (Figure~\ref{fig:rates}).
No numeric feature correlates strongly with stroke (age is highest at 0.25, BMI lowest at 0.04), and no pair of features is strongly correlated (maximum 0.32, age and BMI; Figure~\ref{fig:corr}), so collinearity is not a concern for logistic regression and all features were kept.
All ages were also kept to avoid biasing the model. Most exploration figures are in the Appendix.

\textbf{Model pipeline.}
Figure~\ref{fig:pipeline} shows the full pipeline.
The data was split 80/20 with stratification (\texttt{random\_state=8}), giving 4,087 training and 1,022 test patients with the same stroke share (4.87\% and 4.89\%).
Inside the sklearn Pipeline, numeric features are median-imputed and standardised, and categorical features are one-hot encoded, giving 20 features.
Because the pipeline is refitted inside every cross-validation fold, the median and scaling are learned from training rows only, which prevents data leakage.
Logistic regression with \texttt{class\_weight='balanced'} was tuned with 5-fold stratified cross-validation over seven values of $C$ from $10^{-3}$ to $10^{3}$, scored on F1.
The decision threshold was chosen from cross-validated probabilities on the training rows, and no decisions were made on the test set, which was used only for the final evaluation.
